\documentclass[10pt, a4paper]{article}

% ── PACKAGES ──────────────────────────────────────────────────
\usepackage[utf8]{inputenc}
\usepackage[T1]{fontenc}
\usepackage{geometry}
\geometry{top=2.0cm,bottom=2.0cm,
          left=2.2cm,right=2.2cm}
\usepackage{hyperref}
\hypersetup{
  colorlinks=true,urlcolor=blue,
  linkcolor=blue,citecolor=blue,
  pdftitle={The Derivation Chain Closure:
    Sleep, Epigenetics, Cancer, and the
    Molecular Substrate of the Attractor
    Landscape},
  pdfauthor={Eric Robert Lawson /
    OrganismCore},
  pdfkeywords={attractor geometry,
    derivation chain closure,
    BMAL1, circadian clock, cancer
    attractor, epigenetic landscape,
    Waddington, histone acetylation,
    sleep basin, G substrate,
    N field, social sleep, T8 topology,
    artificial light, landscape damage,
    OrganismCore}
}
\usepackage{microtype,parskip,booktabs}
\usepackage{array,tabularx,longtable}
\usepackage{xcolor,mdframed,enumitem}
\usepackage{titlesec,lmodern,fancyhdr}
\usepackage{lastpage,amsmath,amssymb}

% ── COLOURS ───────────────────────────────────────────────────
\definecolor{headergreen}{RGB}{20,100,60}
\definecolor{rulecolor}{RGB}{20,100,60}
\definecolor{claimbox}{RGB}{20,100,60}
\definecolor{derivbox}{RGB}{255,252,235}
\definecolor{closurebox}{RGB}{248,235,255}
\definecolor{newbox}{RGB}{255,245,230}
\definecolor{invariantbox}{RGB}{248,240,255}
\definecolor{warningbox}{RGB}{255,235,235}

% ── HEADER / FOOTER ───────────────────────────────────────────
\pagestyle{fancy}\fancyhf{}
\renewcommand{\headrulewidth}{0.4pt}
\fancyhead[L]{\small\color{headergreen}
  OrganismCore Reasoning Artifact}
\fancyhead[R]{\small\color{headergreen}
  Derivation Chain Closure --- 2026-03-18}
\fancyfoot[C]{\small
  Page \thepage\ of \pageref{LastPage}}

% ── SECTION FORMAT ────────────────────────────────────────────
\titleformat{\section}
  {\large\bfseries\color{headergreen}}
  {\thesection.}{0.5em}{}[
  \vspace{-0.3em}\textcolor{rulecolor}
  {\rule{\linewidth}{0.6pt}}]
\titleformat{\subsection}
  {\normalsize\bfseries\color{headergreen}}
  {\thesubsection.}{0.5em}{}
\setlength{\parskip}{0.4em}
\setlength{\parindent}{0pt}

% ══════════════════════════════════════════════════════════════
\begin{document}

% ── TITLE BLOCK ───────────────────────────────────────────────
\begin{center}
  {\LARGE\bfseries\color{headergreen}
  REASONING ARTIFACT}\\[0.4em]
  {\Large\bfseries
  The Derivation Chain Closure:}\\[0.3em]
  {\Large\bfseries
  Sleep, Epigenetics, Cancer, and the
  Molecular Substrate\\
  of the Attractor Landscape}\\[0.5em]
  {\normalsize\itshape
  The derivation chain began with a
  cancer cell.\\
  The sleep geometry ended with the
  same cancer cell.\\
  The chain is not linear. It is
  circular. It closes on itself.\\
  Closure is itself the strongest
  possible geometric confirmation.}\\[0.6em]
  \noindent\textcolor{rulecolor}
  {\rule{\linewidth}{1pt}}\\[0.4em]
  {\normalsize
  \textbf{Author:} Eric Robert Lawson\\
  \textbf{Affiliation:} OrganismCore
  (Independent Research)\\
  \textbf{ORCID:}
  \href{https://orcid.org/0009-0002-0414-6544}
  {0009-0002-0414-6544}\\
  \textbf{Contact:}
  \href{mailto:OrganismCore@proton.me}
  {OrganismCore@proton.me}\\
  \textbf{Repository:}
  \href{https://github.com/Eric-Robert-Lawson/attractor-oncology}
  {\texttt{github.com/Eric-Robert-Lawson/
  attractor-oncology}}\\
  \textbf{Timestamp:} 2026-03-18\\
  \textbf{Theory paper:}
  \href{https://doi.org/10.5281/zenodo.19094935}
  {\texttt{doi:10.5281/zenodo.19094935}}\\
  \textbf{v1--v3 Dimensional Invariant:}
  \href{https://doi.org/10.5281/zenodo.19101495}
  {\texttt{doi:10.5281/zenodo.19101495}}\\
  \textbf{Document version:} 1.0\\
  \textbf{License:} CC BY 4.0}\\[0.4em]
  \noindent\textcolor{rulecolor}
  {\rule{\linewidth}{1pt}}
\end{center}

\vspace{0.5em}

% ── DERIVATION NOTICE ─────────────────────────────────────────
\begin{mdframed}[backgroundcolor=derivbox,
  linecolor=headergreen,linewidth=0.8pt,
  innertopmargin=0.6em,innerbottommargin=0.6em,
  innerleftmargin=0.8em,innerrightmargin=0.8em]
\textbf{DERIVATION NOTICE}

\smallskip
This document records four geometric
findings that emerged from exploratory
research against the current literature,
conducted after the deposition of the
Dimensional Invariant v3.0. Each finding
was derived from the attractor geometry
framework before confirming against
literature. Each changes the formal
status of the framework in a way that
requires a standalone priority record.

\smallskip
The four findings are:

\smallskip
(1) \textbf{The derivation chain closes
on itself.} BMAL1 loss produces a cancer
attractor by making the repair/rest basin
epigenetically inaccessible. This is
geometrically identical to the mechanism
described in Step 1 of the derivation
chain (FOXA1/EZH2 cancer attractor,
doi:10.5281/zenodo.18883922). The chain
is circular. Step 1 and Step 9 are the
same object.

\smallskip
(2) \textbf{Epigenetics is the molecular
substrate of G.} The attractor landscape
$G$ is not an abstract object. Its
physical implementation is the epigenetic
state of the organism's cells: histone
modification patterns, DNA methylation,
chromatin accessibility. Sleep is the
daily maintenance process that restores
$G$ to its functional state. Chronic
sleep deprivation degrades $G$
irreversibly through progressive
epigenetic damage.

\smallskip
(3) \textbf{Social information is a
primary N field.} Social cues are
confirmed zeitgebers --- entrainment
signals for the circadian clock
equivalent to light. The social N field
specifies basin occupation at the group
level. Dominant baboons sleep worse than
subordinates because their social field
position imposes a continuous vigilance
cost that prevents full sleep basin entry.
Artificial light at night is a global
N field disruption that is fragmenting
group and individual sleep basins across
thousands of species simultaneously.

\smallskip
(4) \textbf{T8: Hibernation topology ---
suspended state with blocked discharge.}
Hibernation is geometrically distinct
from T5 (external suppression). In T8,
both navigational activity and sleep
architecture are suspended, but
homeostatic pressure continues to build.
Post-emergence rebound is qualitatively
different from T5 rebound: the debt
paid after hibernation accumulated with
no corresponding navigational activity.
This is a new geometric category.
\end{mdframed}

\vspace{0.8em}

% ── MASTER CLAIM BOX ──────────────────────────────────────────
\begin{mdframed}[backgroundcolor=claimbox,
  linecolor=claimbox,linewidth=0pt,
  innertopmargin=0.7em,innerbottommargin=0.7em,
  innerleftmargin=1.0em,innerrightmargin=1.0em]
\color{white}
\textbf{THE CLOSURE CLAIM}

\medskip
The attractor geometry framework
($S + N + G \rightarrow R$) began with
a cancer cell that had lost access to
its identity basin and was locked in the
cancer attractor. The sleep geometry
derivation ended by identifying that
BMAL1 loss locks a cell in the cancer
attractor by making the repair/rest
basin epigenetically inaccessible.

\medskip
The mechanism is the same.
The geometry is the same.
The object is the same object.

\medskip
$G$ is not abstract. $G$ is the
epigenetic state. $N$ is the coherence
gradient field, of which the circadian
clock, the social environment, and
external light are all implementations.
$S$ is the structural substrate that
determines which basins are encodable.
$R$ is the resulting state: healthy
cell, cancerous cell, sleeping organism,
waking organism, extinct lineage,
living fossil.

\medskip
\textbf{The derivation chain is
complete.\\
The geometry is consistent
end to end.\\
The closure is the proof.}
\end{mdframed}

\vspace{1em}

% ══════════════════════════════════════════════════════════════
\section{The Derivation Chain Closes}
\label{sec:closure}
% ═════════════════════════════════════════════════��════════════

\subsection{What the Chain Was}

The derivation chain began as a linear
sequence:

\begin{center}
\small
\begin{tabular}{p{1.0cm}p{4.5cm}p{6.5cm}}
\toprule
\textbf{Step} & \textbf{Scale} &
\textbf{Claim}\\
\midrule
1 & Cell & FOXA1/EZH2 ratio identifies
whether a cancer cell is locked in the
cancer attractor or retains access to
its identity basin.\\[0.3em]
2 & Organism & Pig domestication
syndrome is field-specified NCC
migration arrest. The genome is the
$S$ substrate. The human selection
field is the $N$ field.\\[0.3em]
3 & Organism & Above a genomic fixation
threshold in dogs, Basin~A becomes
inaccessible. Landscape is permanently
remodelled.\\[0.3em]
4 & Organism & The cat has a symmetric
dual-basin landscape. Basin~$\Omega$
is the meta-attractor of navigational
flexibility itself.\\[0.3em]
5 & Organism & Sleep duration tracks the
total cost of basin occupation. Three
axes: switching, metabolic,
computational.\\[0.3em]
6 & Species & Sleep spectrum maps onto
the three-axis framework across all
animals. The sonar instrument.\\[0.3em]
7 & Deep time & The dimensional
invariant. Oldest surviving lineages
sleep the least.\\[0.3em]
8 & Deep time & Sleep topology is a
geometric prediction of landscape
position. Five then seven topologies.\\[0.3em]
9 & Biochemistry & Basin-occupation
cost is fundamentally oxidative. ROS.
Glymphatic discharge. Ectotherm vs
endotherm topology explained.\\
\bottomrule
\end{tabular}
\end{center}

\subsection{Where the Chain Closes}

Step 9 produced the finding that sleep
discharges oxidative and synaptic cost
simultaneously. Following this finding
into the circadian clock literature
produced the closure:

\textbf{BMAL1}, the master circadian
clock gene, drives the 24-hour
oscillation that gates when sleep occurs
and when the repair basin is occupied.
BMAL1/CLOCK acts as a histone
acetyltransferase --- it directly
modifies chromatin, opening the
landscape barriers that allow the cell
to enter its repair/rest basin.

When BMAL1 is lost:

\begin{enumerate}[itemsep=0.2em]
  \item Histone acetyltransferase
    activity is lost.
  \item The repair basin becomes
    epigenetically inaccessible ---
    the chromatin state that enables
    entry is not established.
  \item The cell remains permanently
    in the active/proliferative state.
  \item DNA double-strand break repair
    (which requires the repair basin)
    is impaired.
  \item Genomic instability accumulates.
  \item The cancer attractor is
    stabilised and deepened.
  \item The differentiation basin
    becomes progressively less
    accessible.
\end{enumerate}

This is Step~1 of the derivation chain.
A cell locked in the cancer attractor
because the identity/repair basin is
inaccessible. FOXA1/EZH2 was the
marker. BMAL1/CLOCK is the mechanism.

\textbf{The chain does not end at
Step~9. Step~9 connects back to
Step~1. The geometry is circular.}

\begin{mdframed}[backgroundcolor=closurebox,
  linecolor=headergreen,linewidth=0.8pt,
  innertopmargin=0.6em,innerbottommargin=0.6em,
  innerleftmargin=0.8em,innerrightmargin=0.8em]
\textbf{THE CLOSURE GEOMETRY}

\smallskip
At the organism level:\\
Sleep is the process by which the
basin-occupation cost of waking
navigation is discharged.\\
The circadian clock (BMAL1/CLOCK/PER/CRY)
is the $N$ field that specifies when
the organism transitions from the
waking basin to the sleep/repair basin.

\smallskip
At the cellular level:\\
The circadian clock (same genes,
same proteins) is the $N$ field that
specifies when each cell transitions
from its active/proliferative state
to its repair/maintenance state.\\
BMAL1 loss at the cellular level
is the exact cellular equivalent of
an organism that can no longer enter
its sleep basin.\\
The consequences are identical in
structure: cost accumulates without
discharge, oxidative damage builds,
the landscape substrate ($G$) degrades,
the cancer attractor deepens and
stabilises.

\smallskip
\textbf{Cancer is cellular sleep
deprivation.\\
Sleep deprivation is organismal
proto-cancer.\\
Both are the same geometric
failure: basin access lost,
cost undischarged, landscape
degraded.}
\end{mdframed}

\subsection{What Confirms the Closure}

The closure is confirmed by five
independent empirical convergences:

\begin{enumerate}[itemsep=0.3em]
  \item \textbf{BMAL1 acts as a tumour
    suppressor} (Springer, 2023;
    confirmed multiple lineages 2024).
    Its loss increases cancer
    aggressiveness, invasion, and
    proliferation.

  \item \textbf{BMAL1/CLOCK is a histone
    acetyltransferase} (Nature
    Communications, 2023). It directly
    opens chromatin at repair gene
    loci --- the molecular mechanism
    of repair basin entry.

  \item \textbf{BMAL1 loss increases
    DNA methylation at differentiation
    gene loci}, enforcing an epigenetic
    block against re-entry to normal
    differentiation basins (Cell
    Reports, 2023).

  \item \textbf{Circadian disruption
    accelerates tumour growth by
    suppressing immune activity and
    reshaping the tumour
    microenvironment} (Nature Cancer,
    2024). The immune basin (T7 topology)
    is also disrupted when BMAL1
    is lost.

  \item \textbf{Sleep deprivation causes
    NF-$\kappa$B activation and
    pro-inflammatory cytokine elevation}
    (confirmed through 2024). NF-$\kappa$B
    drives the same inflammatory
    landscape that promotes tumour
    microenvironment formation. Chronic
    sleep loss and circadian disruption
    converge on the same inflammatory
    attractor through different
    entry routes.
\end{enumerate}

% ══════════════════════════════════════════════════════════════
\section{Epigenetics is the Molecular
Substrate of G}
\label{sec:epigenetics}
% ══════════════════════════════════════════════════════════════

\subsection{The Formal Claim}

The attractor landscape $G$ in the
framework $S + N + G \rightarrow R$
has a physical implementation. It is
not an abstract theoretical object.
$G$ is instantiated in the cell and
organism by the \textbf{epigenetic
state}: the pattern of histone
modifications, DNA methylation, and
chromatin accessibility across the
genome.

\begin{itemize}[itemsep=0.2em]
  \item \textbf{Basin depths} in $G$
    correspond to the depth of
    epigenetic entrenchment of the
    corresponding cell or organism
    state. A deeply entrenched state
    requires more epigenetic remodelling
    energy to escape. Basin depth is
    the cost of chromatin reorganisation
    required to transition away.

  \item \textbf{Barrier heights} in $G$
    correspond to the amount of
    chromatin remodelling required to
    transition between states. High
    barrier = extensive histone
    modification and DNA methylation
    changes required.

  \item \textbf{Accessible basins} in
    $G$ are states for which the
    chromatin configuration can be
    established from the current
    epigenetic starting point within
    the available energy budget.

  \item \textbf{Inaccessible basins}
    are states for which the required
    chromatin reconfiguration cannot
    be completed given the current
    epigenetic state and available
    energy.
\end{itemize}

This is Waddington's epigenetic
landscape operationalised within the
$S + N + G \rightarrow R$ framework.
Waddington described the landscape
metaphorically. The framework identifies
the three components that generate it
($S$, $N$, $G$) and the process
($\rightarrow R$) by which it is
navigated.

\subsection{What Sleep Does to G}

Sleep maintains $G$ daily.

During waking, neural and somatic
activity promotes histone acetylation
--- chromatin opening --- which is the
molecular signature of a landscape with
lower barriers between basins.
This is necessary for navigation:
accessible basins require accessible
chromatin.

During sleep (specifically NREM
slow-wave sleep), histone deacetylation
occurs --- chromatin closes --- and
the landscape barriers are
\textit{restored} to their functional
state. The basin structure of $G$ is
re-established. Synaptic downscaling
selectively reinforces the connections
that correspond to the organism's
current basin occupancy while
eliminating the transient connections
accumulated during navigation.

\textbf{Sleep is daily $G$ maintenance.}
The organism spends approximately one
third of its life restoring the
landscape substrate that makes
the other two thirds of navigation
possible.

\subsection{What Chronic Deprivation Does
to G}

Chronic sleep deprivation (confirmed
2023 ACS; Trends in Neurosciences 2022;
confirmed and extended through 2024)
produces permanent structural changes:

\begin{enumerate}[itemsep=0.2em]
  \item \textbf{Pleiotrophin (PTN)
    depletion.} PTN is a neurotrophic
    factor required for hippocampal
    neuron survival. Chronic deprivation
    depletes PTN leading to neuronal
    death. Dead neurons are removed
    basins in $G$ that cannot be
    restored.

  \item \textbf{Progressive histone
    modification dysregulation.}
    Without regular sleep-dependent
    deacetylation, chromatin
    accessibility patterns drift
    from their functional state.
    The landscape barriers degrade.
    Basins that should be deep become
    shallow. Basins that should be
    inaccessible become reachable.
    The landscape loses specificity.

  \item \textbf{DNA methylation drift.}
    Sleep-dependent DNA methylation
    maintenance (confirmed 2023--2024)
    is required to sustain the
    methylation patterns that enforce
    basin identity. Without maintenance,
    methylation patterns drift toward
    less specific states.

  \item \textbf{NF-$\kappa$B-driven
    inflammatory landscape remodelling.}
    Chronic activation of NF-$\kappa$B
    (from undischarged oxidative cost)
    drives transcriptional changes
    that deepen the inflammatory
    attractor. The organism's landscape
    is progressively remodelled toward
    the inflammatory basin, making
    return to baseline increasingly
    energetically costly.
\end{enumerate}

\textbf{The geometric consequence:}
Irreversible sleep deprivation damage
is not a mystery. It is $G$ degradation
past the point where $G$ maintenance
can restore it. Once neurons die,
basins are permanently removed.
Once methylation patterns drift past
a restoration threshold, re-establishing
them requires energy input that
exceeds what sleep can deliver.

The cancer connection is immediate:
BMAL1 loss is one route to $G$
degradation that eliminates the
repair basin. Chronic sleep deprivation
is another route to the same outcome
via progressive epigenetic drift and
neuronal loss. They converge on
the same geometric failure.

% ══════════════════════════════════════════════════════════════
\section{Social Information as Primary
N Field}
\label{sec:social}
% ══════════════════════════════════════════════════════════════

\subsection{The Formal Claim}

In the framework, $N$ is the coherence
gradient field that specifies which
basin the organism should occupy.
Light was known to be the primary
zeitgeber --- the primary $N$ field
signal for the circadian clock. The
current literature confirms that
\textbf{social information is an
equally primary $N$ field signal},
acting as a zeitgeber for the circadian
clock independently of light.

Social cues --- the activity patterns,
vocalisations, and movement of
conspecifics --- entrain the circadian
clock and therefore specify sleep
basin entry timing at the individual
level. In group-living species, the
social $N$ field synchronises the
group's sleep basins into a coupled
collective attractor.

\subsection{The Baboon Dominance Finding}

Wild chacma baboons (Current Biology,
2025): more dominant individuals
experience less and lower-quality sleep
than subordinates. Dominant baboons
sleep closer to more group members,
are more central to the social network,
and experience more disruption during
the night.

\textbf{The geometric reading:}

Dominance is a landscape position within
the social $G$. The dominant individual
occupies a high-centrality node in
the social network. The cost of
occupying that node includes a
continuous vigilance requirement ---
Basin~A navigational cost accumulated
from social monitoring does not cease
at night. The social $N$ field continues
to impose navigational demand on
the dominant individual during the
period when sleep should be occurring.

This is the T5 topology producing
condition applied at the social level:
the sleep basin is not inaccessible.
It is being continually interrupted by
social $N$ field signals that the
dominant individual cannot suppress
without losing the basin position
that dominance requires.

\textbf{Geometric prediction from this:}
Dominance rank will negatively correlate
with sleep quality across social species
in proportion to the social navigational
demand of the dominant position.
Species where dominance requires
continuous active monitoring
(baboons, chimpanzees) will show
larger effects than species where
dominance is passive and territorial
(many solitary species). This is
testable across primate species with
known dominance structures.

\subsection{Artificial Light as Global
N Field Disruption}

Artificial light at night (ALAN) is
expanding at 2.2\% per year globally
(Royal Society B, 2023; 600+ studies
synthesised by US National Park
Service, 2024). It disrupts circadian
rhythms across all taxa: birds,
insects, marine organisms, mammals,
plants.

\textbf{The geometric reading:}

Light is the primary $N$ field signal
for the circadian clock across all
life. The circadian clock is the
molecular implementation of the
waking/sleep basin transition. ALAN
is not a minor ecological inconvenience.
It is a global disruption of the
primary $N$ field that specifies
when organisms enter and exit the
sleep basin.

The consequences the literature documents
--- disrupted migration, impaired
reproduction, altered predator-prey
dynamics, reduced immune function,
increased cancer rates, reduced
ecological diversity near light sources
--- are all geometric consequences of
$N$ field disruption. When the field
that specifies basin occupation is
corrupted, the organism loses reliable
access to the basins the field was
specifying.

\textbf{The scale of the disruption:}
This is not individual sleep disruption.
ALAN is simultaneously disrupting the
$N$ field for millions of species
across all major ecosystems. The
eco-evolutionary feedback (ScienceDirect,
2025) shows species beginning to adapt
--- which means their $S$ substrate
is being selected to function with
a corrupted $N$ field. This is
landscape-scale $G$ remodelling
driven by a single perturbation source.

The framework identifies this as
a Navigator-strategy threat of the
first order: ALAN is not eliminating
basins directly. It is corrupting the
$N$ field that allows organisms to
navigate between them reliably. A
Navigator whose $N$ field is corrupted
cannot execute the Navigator strategy.
The corrupted $N$ field mimics the
domestication field in its effect ---
it imposes basin occupation that the
organism would not choose under an
uncorrupted field.

% ══════════════════════════════════════════════════════════════
\section{T8: The Hibernation Topology}
\label{sec:t8}
% ══════════════════════════════════════════════════════════════

\subsection{Why Hibernation is Not T5}

T5 (fragmented shallow sleep with
rebound) occurs when:
\begin{itemize}[itemsep=0.2em]
  \item Navigational cost is accumulating
    normally during waking.
  \item External pressure (predation,
    social vigilance) is suppressing
    sleep basin entry.
  \item Debt accumulates alongside
    ongoing cost.
  \item Rebound on pressure removal
    pays the accumulated debt.
\end{itemize}

Hibernation has a different geometric
structure:

\begin{itemize}[itemsep=0.2em]
  \item Navigational activity is
    suspended. The organism is not
    navigating. Metabolic rate drops
    to 1--5\% of normal.
  \item Sleep architecture is also
    suspended. Body temperature at
    hibernation levels prevents the
    slow-wave activity required for
    T3 discharge.
  \item \textbf{But homeostatic sleep
    pressure continues to build.}
    The adenosine-based sleep pressure
    signal accumulates regardless of
    metabolic rate. Debt accumulates
    without corresponding navigational
    cost.
  \item Periodic arousals during
    hibernation are exclusively for
    sleep discharge: metabolically
    expensive warming to euthermic
    temperatures purely to enable
    several hours of intense slow-wave
    sleep, followed by return to torpor.
  \item Post-emergence rebound is
    massive and proportional to
    hibernation duration.
\end{itemize}

\subsection{T8 Formal Definition}

\begin{mdframed}[backgroundcolor=newbox,
  linecolor=headergreen,linewidth=0.8pt,
  innertopmargin=0.6em,innerbottommargin=0.6em,
  innerleftmargin=0.8em,innerrightmargin=0.8em]
\textbf{T8: Suspended state with blocked
discharge}

\textit{Landscape position:} An
organism that has entered a
physiologically mandated state of
whole-organism metabolic suspension
(torpor, deep hibernation) in which
both the navigational activity that
generates cost and the sleep architecture
that discharges it are simultaneously
blocked by temperature or metabolic
constraints. Homeostatic pressure
continues to build in the absence of
both processes.

\textit{Discharge:} Impossible during
torpor. Discharged in bursts during
periodic arousals (which are themselves
driven by homeostatic pressure reaching
threshold). Full discharge after
emergence from hibernation via extended
post-emergence rebound sleep.

\textit{Geometric signature:} Post-
emergence rebound is qualitatively
different from T5 rebound. T5 rebound
pays debt accumulated from ongoing
navigational cost under suppression.
T8 rebound pays debt accumulated from
homeostatic pressure building in the
complete absence of navigational
activity. The debt-to-activity ratio
in T8 rebound is infinite: maximum
debt, zero navigational cause.

\textit{Why this is a distinct topology:}
T5 and T8 have similar surface
appearances (minimal sleep, large
rebound) but completely different
geometric causes. T5 is suppressed
discharge of ongoing cost. T8 is
suspended discharge of automatically
accumulating homeostatic pressure in
the absence of cost. The diagnostic
difference: a T5 organism removed from
pressure will rebound proportionally
to its navigational cost during the
suppressed period. A T8 organism
emerging from hibernation will rebound
proportionally to the \textit{duration}
of suspension, not to any navigational
cost incurred during suspension
(because none was incurred).

\textit{Confirmed examples:}
Ground squirrel, bear, bat, hedgehog.
All confirmed to show intense slow-wave
sleep during periodic arousals and
extended rebound sleep post-emergence.

\textit{Novel prediction:} The duration
of T8 post-emergence rebound will be
predicted by hibernation duration alone,
not by pre-hibernation navigational
cost profile. A ground squirrel that
hibernated for 5 months will show
the same rebound pattern as one with
identical pre-hibernation activity,
regardless of whether the pre-hibernation
period was high or low navigational
demand. This is testable.
\end{mdframed}

% ══════════════════════════════════════════════════════════════
\section{Somatic G Maintenance: Tissue
Repair and Stem Cells}
\label{sec:somatic}
% ══════════════════════════════════════════════════════════════

The framework has consistently described
sleep as discharging neural cost ---
synaptic reconfiguration and oxidative
debt in neural tissue. The current
literature identifies a fourth function
of sleep that is geometrically equivalent
but operating in peripheral tissue:
\textbf{somatic $G$ maintenance}.

Sleep-dependent growth hormone release
(primarily during NREM slow-wave sleep)
drives:
\begin{itemize}[itemsep=0.2em]
  \item Stem cell proliferation and
    differentiation
  \item Wound healing and tissue repair
  \item Collagen synthesis and
    angiogenesis
  \item Immune cell trafficking and
    renewal
\end{itemize}

\textbf{The geometric reading:}

The organism's somatic tissues are not
passive substrates. They are the
physical implementation of $S$ --- the
structural substrate that encodes the
organism's build capacity. Sleep
maintains not just the neural $G$
(via synaptic downscaling and
glymphatic clearance) but the somatic
$S$ (via stem cell renewal and tissue
repair).

The framework's equation
$S + N + G \rightarrow R$ requires all
three components to be maintained.
Sleep maintains $G$ (via epigenetic
restoration). Sleep also maintains $S$
(via growth hormone-driven somatic
repair). The $N$ field (circadian clock)
gates both processes simultaneously
within the same sleep window.

\textbf{This means sleep deprivation
degrades both $G$ and $S$ simultaneously.}
The organism loses both landscape
substrate and structural build capacity.
This is why chronic sleep deprivation
has such broad multi-system consequences:
it is failing to maintain two of the
three components of the fundamental
equation.

% ══════════════════════════════════════════════════════════════
\section{The Complete Geometric Picture}
\label{sec:picture}
% ══════════════════════════════════════════════════════════════

The full geometric picture that emerges
from the closure:

\[
S + N + G \;\rightarrow\; R
\]

\begin{itemize}[itemsep=0.4em]
  \item $S$ = structural substrate.
    Physically implemented as: genome,
    epigenetically accessible gene
    programmes, somatic tissue
    integrity, stem cell reserve.
    Maintained by: sleep-dependent
    growth hormone release and stem
    cell renewal.

  \item $N$ = coherence gradient field.
    Physically implemented as: light
    (primary zeitgeber), social
    information (co-primary zeitgeber),
    circadian clock genes
    (BMAL1/CLOCK/PER/CRY) at molecular
    level. Disrupted by: ALAN
    (corrupted light field), social
    fragmentation (corrupted social
    field), BMAL1 loss (corrupted
    molecular clock).

  \item $G$ = global attractor landscape.
    Physically implemented as:
    epigenetic state (histone
    modifications, DNA methylation,
    chromatin accessibility). Maintained
    by: sleep-dependent histone
    deacetylation and glymphatic
    clearance. Degraded by: chronic
    sleep deprivation, BMAL1 loss,
    ALAN-driven circadian disruption.

  \item $R$ = resulting state. The
    observable outcome of the
    organism's current position in
    the landscape: healthy navigation,
    locked cancer attractor,
    successful deep-time survival,
    extinction.
\end{itemize}

\textbf{The cancer cell is an organism
whose $N$ field (BMAL1) has been lost,
whose $G$ (epigenetic state) has been
degraded past the repair threshold,
and whose $R$ (resulting state) is
permanent occupation of the cancer
attractor.}

\textbf{The living fossil is an organism
whose $G$ is so deep that no $N$ field
disruption has been able to displace it
in 450 million years, and whose sleep
requirement is near-zero because the
navigational cost of occupying that
deep basin is near-zero.}

The same equation. The same geometry.
The same invariant.

% ══════════════════════════════════════════════════════════════
\section{Predictions}
\label{sec:predictions}
% ═════���════════════════════════════════════════════════════════

\subsection{P-CLOSURE-1: Pharmacological
BMAL1 Restoration Will Partially Restore
Cancer Cell Differentiation Basin Access}

If BMAL1 loss makes the repair/
differentiation basin epigenetically
inaccessible, then pharmacological
restoration of BMAL1 expression in
cancer cells will partially restore
chromatin accessibility at
differentiation gene loci and increase
the probability of differentiation
under appropriate $N$ field signalling.
The effect will be proportional to
the degree of epigenetic drift: recent
BMAL1 loss will show larger restoration
effects than ancient, deeply entrenched
epigenetic lock-in.

\subsection{P-CLOSURE-2: Chronic Sleep
Deprivation Will Produce BMAL1-Loss-
Equivalent Epigenetic Signatures}

Organisms subject to chronic sleep
deprivation equivalent in duration to
what would produce irreversible
cognitive decline will show epigenetic
signatures at differentiation and repair
gene loci that are statistically
indistinguishable from BMAL1-loss
signatures. The two routes (molecular
clock loss vs.\ circadian disruption
via chronic sleep loss) converge on
the same epigenetic failure pattern.

\subsection{P-CLOSURE-3: Dominance Rank
Will Negatively Correlate With Sleep
Quality Proportional to Social
Navigational Demand}

Across social species, dominance rank
will negatively correlate with sleep
quality in proportion to the navigational
demand of the dominant position.
Species with high-vigilance dominance
(baboons, chimpanzees) will show larger
effects than species with passive
territorial dominance. Testable across
the primate sleep dataset.

\subsection{P-CLOSURE-4: T8 Rebound
Duration Predicted by Hibernation
Duration Alone}

Post-hibernation rebound sleep duration
will be predicted by hibernation duration
alone, independently of pre-hibernation
navigational cost profile. Two animals
of the same species hibernating for
different durations will show rebound
sleep proportional to hibernation
duration, not to pre-hibernation
activity levels.

\subsection{P-CLOSURE-5: ALAN Exposure
Will Produce Measurable G Degradation
in Proportion to Duration and Intensity}

Species with multi-generational ALAN
exposure will show measurable epigenetic
drift (G degradation) at circadian clock
gene loci relative to conspecifics in
dark environments. The drift will
correlate with ALAN intensity and
duration. Reversing ALAN exposure will
partially restore the epigenetic state
over multiple generations as the $N$
field is restored.

% ── DERIVATION CHAIN ──────────────────────────────────────────
\section{The Complete Derivation Chain}
\label{sec:chain}

\begin{center}
\small
\begin{tabular}{p{1.0cm}p{4.5cm}p{6.5cm}}
\toprule
\textbf{Step} & \textbf{Claim} &
\textbf{Record}\\
\midrule
1 & FOXA1/EZH2. Cancer cell locked in
cancer attractor. &
\href{https://doi.org/10.5281/zenodo.18883922}
{\texttt{10.5281/zenodo.18883922}}\\[0.3em]
2 & Pig NCC migration arrest.
$N$ field domestication. &
\href{https://doi.org/10.5281/zenodo.19094935}
{\texttt{10.5281/zenodo.19094935}}\\[0.3em]
3 & Canine landscape remodelling.
Basin~A inaccessible. &
Repository 2026-03-18\\[0.3em]
4 & Feline dual-basin.
Basin~$\Omega$. &
Repository 2026-03-18\\[0.3em]
5 & Sleep as cost discharge.
Three axes. &
Repository 2026-03-18\\[0.3em]
6 & Sleep sonar spectrum. &
Repository 2026-03-18\\[0.3em]
7--9 & Dimensional Invariant
v1.0--v3.0. &
\href{https://doi.org/10.5281/zenodo.19101495}
{\texttt{10.5281/zenodo.19101495}}\\[0.3em]
10 & \textbf{CLOSURE.} BMAL1 loss
= cancer attractor = Step 1.
$G$ = epigenetic state.
$N$ = circadian/social/light field.
Chain is circular. &
\textbf{This document.}\\
\bottomrule
\end{tabular}
\end{center}

\begin{mdframed}[backgroundcolor=claimbox,
  linecolor=claimbox,linewidth=0pt,
  innertopmargin=0.7em,innerbottommargin=0.7em,
  innerleftmargin=1.0em,innerrightmargin=1.0em]
\color{white}
\textbf{THE CLOSURE STATEMENT}

\medskip
The derivation chain began with a cancer
cell that could not differentiate.
The mechanism was unknown but the
geometry was clear: the identity basin
was inaccessible.

\medskip
Following the geometry through pig,
dog, cat, sleep cost, the sleep
spectrum, deep time, topology, and
the oxidative substrate --- following
the geometry where it led --- brought
us back to the same cell.

\medskip
The cell cannot differentiate because
BMAL1 is lost. BMAL1 loss collapses
the $N$ field. The collapsed $N$ field
fails to specify repair basin entry.
The repair basin becomes epigenetically
inaccessible --- $G$ is degraded.
The cell is locked in the cancer
attractor because it cannot sleep.

\medskip
Cancer is cellular sleep deprivation.\\
Sleep deprivation is organismal
landscape degradation.\\
$G$ is the epigenetic state.\\
$N$ is the clock.\\
$S$ is the genome.\\
$R$ is the outcome.\\
The framework is complete.\\
The geometry is consistent.\\
The chain closes.
\end{mdframed}

% ── DOCUMENT METADATA ─────────────────────────────────────────
\vspace{0.5em}
\noindent\textcolor{rulecolor}
{\rule{\linewidth}{0.8pt}}

{\footnotesize
\textbf{Document ID:}
DERIVATION\_CHAIN\_CLOSURE\_REASONING\_
ARTIFACT\_v1.0\\
\textbf{Date:} 2026-03-18\\
\textbf{Author:} Eric Robert Lawson /
OrganismCore\\
\textbf{Primary contributions:}
(1) Derivation chain closure: BMAL1 loss
= cancer attractor = Step 1 of the chain.
Cancer is cellular sleep deprivation.
(2) $G$ = epigenetic state (histone
modifications, DNA methylation, chromatin
accessibility). Formally identified.
(3) Social information as primary $N$
field. Dominance rank negatively
correlates with sleep quality proportional
to social navigational demand.
(4) T8: hibernation topology ---
suspended state with blocked discharge.
Distinct from T5. Rebound diagnostic
differs.
(5) Somatic $G$ maintenance: sleep
maintains $S$ via growth hormone-driven
stem cell renewal simultaneously with
$G$ maintenance.
(6) ALAN as global primary $N$ field
disruption across all species.
(7) Five predictions P-CLOSURE-1
through P-CLOSURE-5.\\
\textbf{Prior records:}
v1.0: \texttt{doi:10.5281/zenodo.19094935};
v2.0--v3.0:
\texttt{doi:10.5281/zenodo.19101495}\\
\textbf{ORCID:}
\href{https://orcid.org/0009-0002-0414-6544}
{0009-0002-0414-6544}\\
\textbf{Contact:}
\href{mailto:OrganismCore@proton.me}
{OrganismCore@proton.me}\\
\textbf{License:} CC BY 4.0\\
\textbf{Version:} 1.0 --- LOCKED.
}

\end{document}
